\documentclass[10pt,twoside,slovak,a4paper]{article}

\usepackage[slovak]{babel}
%\usepackage[T1]{fontenc}
\usepackage[IL2]{fontenc} % lepšia sadzba písmena Ľ než v T1
\usepackage[utf8]{inputenc}
\usepackage{graphicx}
\usepackage{url} % príkaz \url na formátovanie URL
\usepackage{hyperref} % odkazy v texte budú aktívne (pri niektorých triedach dokumentov spôsobuje posun textu)
\usepackage{todonotes}
\usepackage{cite}

\title{Zameranie projektu}
\author{Mykola Kornilov, Danyil Komarchuk, Kholoshyn Kostiantyn, 4\\
	{\small Slovenská technická univerzita v Bratislave}\\
	{\small Fakulta informatiky a informačných technológií}\\
	}
\date{3. Oktobra 2026}
\begin{document}

\maketitle
\section*{\textbf VV-A1-03	
\normalsize Názov projektu/Project title in English}
NoiseMap: mapovanie hluku v interiéroch v reálnom čase pomocou senzorových sietí a crowdsourcingových meraní 
Real-time indoor noise mapping using sensor networks and crowdsourced measurements\\[2pt]

\section*{\textbf VV-A1-04	
\normalsize Akronym projektu/Acronym of the project}
NM\\[2pt]

\section*{\textbf VV-A1-09
\normalsize Anotácia(max. 2000 znakov)}
Keď sa ľudia snažia nájsť v budove tiché miesto na prácu alebo štúdium, často musia stráviť veľa času
jeho hľadaním.\\[2pt]

Projekt NoiseMap sa preto zameriava predovšetkým na vytvorenie platformy, ktorá by zbierala údaje o hluku v rôznych budovách (napr. knižniciach, kaviarňach, coworkingových priestoroch) a prezentovala ich vo forme interaktívnej mapy spolu s typickými dennými vzormi.\\[2pt]

Očakáva sa, že údaje budú pochádzať zo senzorov inštalovaných v budovách, ako aj z dobrovoľných meraní návštevníkov prostredníctvom mobilnej aplikácie, pričom sa zohľadní ochrana súkromia.\\[2pt]

Pozornosť sa bude sústreďovať na spracovanie a overovanie týchto údajov s využitím efektívnych metód na
zachytávanie čistého zvuku.\\[2pt]

Výsledkom bude užívateľsky prívetivá aplikácia, v ktorej sa všetko zobrazí na interaktívnej mape.\\[2pt]

\section*{\textbf VV-A1-09
\normalsize Anotation(max. 2000 characters))}
When people try to find a quiet place in the building to work or study, they often have to spend a lot of time 
looking for it. \\[2pt]

Therefore, the NoiseMap is a project that mainly focuses on creating a platform that would collect noise data in various buildings(e.g. libraries, cafés, coworking spaces) and present it as an interactive map together with typical patterns over the day. \\[2pt]

The data is expected to come from sensors installed in buildings as well as from voluntary measurements by visitors via a mobile app, with privacy protection taken into account. \\[2pt]

Attention will be focused on processing and validating that data supported by using efficient methods to 
capture the pure sound.\\[2pt]

The outcome will be a user-friendly interface application, where everything will be shown in the interactive map.
\end{document}